\section{Conferencia invitada: la filosofía de diseño de Claude Code (Boris Cherny)}

\subsection{La programación está en un punto de inflexión}

Boris Cherny, el creador de Claude Code, compartió una idea clave: la productividad de los lenguajes de programación y de los IDE ha crecido siguiendo una curva exponencial, y la IA es precisamente la fuerza central que impulsa la ola más reciente de ese crecimiento.

\begin{knowledgebox}{El crecimiento exponencial de la productividad en la programación}
De Assembly a FORTRAN, de C a Python, de VS Code a Cursor/Claude Code: cada generación de lenguajes de programación e IDE ha logrado un incremento lineal de la productividad en una escala logarítmica. Las técnicas de verificación (verification) también han evolucionado en paralelo: de la depuración manual a los tipos estáticos, de las pruebas automatizadas al fuzz testing y al self-play impulsados por IA.
\end{knowledgebox}

\subsection{Los principios de diseño de Claude Code}

La filosofía de diseño central de Claude Code se resume en tres «omnipresencias»:

\begin{enumerate}
    \item \textbf{Terminal-native}: se ejecuta en la terminal, no como un complemento de IDE
    \item \textbf{Low-level model access}: proporciona acceso directo al modelo de bajo nivel
    \item \textbf{Infinitely hackable}: personalizable sin límites
\end{enumerate}

\subsection{Un solo Claude Code, múltiples formas}

Claude Code atiende distintos escenarios a través de múltiples formas:

\begin{itemize}
    \item \textbf{Terminal}: uso directo desde la línea de comandos
    \item \textbf{IDE}: integración en editores como VS Code
    \item \textbf{Web \& iOS}: acceso a través de \url{claude.ai/code}
    \item \textbf{GitHub App}: instalación mediante \texttt{/install-github-app}
    \item \textbf{SDK}: invocación programática, con salida en JSON y composición mediante tuberías
\end{itemize}

\subsection{Patrones típicos de uso de Claude Code}

\subsubsection{Preguntas y exploración sobre la base de código}

\begin{itemize}
    \item «How do I make a new ValidationTemplateFactory?»
    \item «Why did we fix issue \#18363 by adding the if/else in login.ts?»
    \item «Look at PR \#9383, then verify which app versions were impacted»
    \item «What did I ship last week?»
\end{itemize}

\subsubsection{Adaptar el flujo de trabajo a la tarea}

Distintas tareas se adaptan mejor a distintos patrones de flujo de trabajo:
\begin{itemize}
    \item \textbf{explore $\rightarrow$ plan $\rightarrow$ confirm $\rightarrow$ code $\rightarrow$ commit}: adecuado para bug fixes que requieren investigación previa
    \item \textbf{tests $\rightarrow$ commit $\rightarrow$ code $\rightarrow$ iterate $\rightarrow$ commit}: patrón TDD (desarrollo guiado por pruebas)
    \item \textbf{code $\rightarrow$ screenshot $\rightarrow$ iterate}: patrón de iteración visual para el desarrollo de UI
\end{itemize}

\subsubsection{Iteración rápida de prototipos}

Boris mostró un ejemplo vívido de iteración de prototipos: mediante instrucciones sucesivas en lenguaje natural, iteró rápidamente el diseño de la UI de una lista de tareas pendientes (TODO); pasando de una vista inicial del uso de herramientas, a un título en línea, a un panel independiente y, finalmente, a una visualización combinada con el spinner, cada iteración requirió solo un prompt.

\begin{importantbox}{Construir para el modelo de dentro de 6 meses}
La idea central que Boris subrayó: no diseñes tu flujo de trabajo para las capacidades del modelo de hoy, sino prepárate para las capacidades del modelo de \textbf{dentro de 6 meses}. Los modelos evolucionan rápidamente, y tus herramientas y flujos de trabajo también deben tener capacidad de evolucionar.
\end{importantbox}

\subsection{Restricciones organizacionales al gestionar múltiples Agents}

Otra conclusión implícita de la Week 4 es que, cuando los desarrolladores empiezan a gestionar varios Agents a la vez, la forma misma de gestionar proyectos también debe cambiar.

\begin{itemize}
    \item \textbf{Las tareas deben poder dividirse}: si los requisitos son ambiguos, es difícil delegarlos con seguridad a varios Agents en paralelo
    \item \textbf{Los límites deben ser claros}: es necesario especificar de antemano qué archivos, qué pasos de verificación y qué granularidad de commits corresponden a cada Agent
    \item \textbf{La revisión debe estandarizarse}: cuanto más se dependa de la ejecución paralela de Agents, más se necesita un mecanismo consistente de revisión de diffs y de reversión
\end{itemize}

\begin{warningbox}{El principal modo de falla del Agent Manager no es que el modelo no sea lo bastante inteligente, sino una división inadecuada de las tareas}
Muchos fracasos en la ejecución paralela no ocurren porque el Agent no pueda hacerlo, sino porque el humano no definió, antes de empezar, límites de responsabilidad ni criterios de aceptación claros. Dicho de otro modo, ser Agent Manager es, ante todo, una capacidad de gestión de proyectos, y solo en segundo lugar, una técnica de prompting.
\end{warningbox}

\subsection{Resumen de la sección}

La filosofía de diseño de Claude Code es terminal-native, acceso a modelos de bajo nivel y personalización sin límites. Cubre todo el SDLC, desde la exploración de la base de código hasta la gestión de CI/CD. La estrategia central de uso consiste en elegir el patrón de flujo de trabajo adecuado según el tipo de tarea, y prepararse con antelación para las capacidades futuras de los modelos.

